\documentclass{beamer}
\usetheme{Madrid}
\usecolortheme{default}
\setbeamertemplate{navigation symbols}{}
\setbeamertemplate{footline}[frame number]


% Reduce overall font size
\renewcommand{\normalsize}{\fontsize{9}{11}\selectfont}
\renewcommand{\small}{\fontsize{8}{10}\selectfont}
\renewcommand{\footnotesize}{\fontsize{7}{9}\selectfont}
\renewcommand{\scriptsize}{\fontsize{6}{8}\selectfont}

% Packages
\usepackage{graphicx}
\usepackage{amsmath, amssymb}
\usepackage{booktabs}
\usepackage{listings}
\usepackage{xcolor}
\usepackage{tikz}

% Define colors
\definecolor{acamsblue}{RGB}{0, 51, 102}
\definecolor{acamsgold}{RGB}{198, 156, 67}
\setbeamercolor{title}{fg=acamsblue}
\setbeamercolor{frametitle}{fg=acamsblue}
\setbeamercolor{structure}{fg=acamsblue}

% Title Information
\title{ACAMS CAMS Exam Mastery Guide}
\subtitle{Advanced AML/CFT Concepts \& Regulatory Frameworks}
\author{Compliance Professional Training}
\date{\today}

\begin{document}
	
	% ============================================================
	% SLIDE 1: TITLE
	% ============================================================
	\begin{frame}
		\titlepage
		\begin{center}
			\small \textit{50+ Slides of Hard-Hitting Exam Content}
		\end{center}
	\end{frame}
	
	% ============================================================
	% SLIDE 2: GLOBAL FRAMEWORK - DETAILED
	% ============================================================
	\begin{frame}{The AML/CFT Ecosystem}
		\begin{block}{Anti-Money Laundering (AML)}
			\begin{itemize}
				\item Laws, regulations, and procedures to prevent criminals from disguising illegally obtained funds as legitimate income.
				\item \textbf{Three Stages:} Placement, Layering, Integration.
				\item \textbf{Predicate Offenses:} Drug trafficking, fraud, corruption, tax evasion, terrorism financing, and over 200 other crimes.
			\end{itemize}
		\end{block}
		\begin{block}{Counter-Financing of Terrorism (CFT)}
			\begin{itemize}
				\item Efforts to prevent the financial support of terrorist organizations and activities.
				\item \textbf{Key Difference from AML:} The source of funds may be legitimate; the \textit{destination} or \textit{purpose} is illicit.
				\item \textbf{FATF Special Recommendations:} Focus on wire transfers, alternative remittance systems, and non-profit organizations.
			\end{itemize}
		\end{block}
		\begin{alertblock}{Critical Concept}
			The global AML/CFT framework is \textbf{risk-based}, not prescriptive. Institutions must tailor controls to their specific risk profiles.
		\end{alertblock}
	\end{frame}
	
	% ============================================================
	% SLIDE 3: FATF - DEEP DIVE
	% ============================================================
	\begin{frame}{Financial Action Task Force (FATF)}
		\begin{columns}[T]
			\column{0.5\textwidth}
			\begin{alertblock}{FATF's Core Functions}
				\begin{itemize}
					\item Sets \textbf{40 Recommendations} (revised 2012, 2019, 2023).
					\item Conducts \textbf{Mutual Evaluations} (assessments of member countries).
					\item Issues \textbf{Public Statements} on high-risk jurisdictions.
					\item Monitors \textbf{emerging risks} (e.g., virtual assets, environmental crime).
				\end{itemize}
			\end{alertblock}
			\column{0.5\textwidth}
			\begin{exampleblock}{FATF Lists}
				\begin{itemize}
					\item \textbf{Black List (Call for Action):} Jurisdictions with strategic AML/CFT deficiencies (e.g., Iran, North Korea).
					\item \textbf{Grey List (Increased Monitoring):} Jurisdictions committed to resolving deficiencies (e.g., Nigeria, South Africa, UAE).
					\item \textbf{Consequence:} Enhanced due diligence (EDD) required for transactions involving black/grey-listed countries.
				\end{itemize}
			\end{exampleblock}
		\end{columns}
		\begin{block}{Exam Tip}
			Know the difference between "black list" and "grey list" and the corresponding obligations for financial institutions.
		\end{block}
	\end{frame}
	
	% ============================================================
	% SLIDE 4: FATF IMMEDIATE OUTCOMES - ADVANCED
	% ============================================================
	\begin{frame}{FATF Immediate Outcomes (IOs)}
		\begin{definition}[Immediate Outcome 3 (IO.3)]
			\textbf{Financial supervision} is risk-based and effective. Supervisors must:
			\begin{itemize}
				\item Use a \textbf{risk-based approach} to allocate supervisory resources.
				\item Conduct \textbf{on-site and off-site} inspections.
				\item Apply \textbf{proportionate sanctions} for non-compliance.
			\end{itemize}
		\end{definition}
		\begin{definition}[Immediate Outcome 4 (IO.4)]
			\textbf{Financial institutions and DNFBPs} apply preventive measures consistent with their ML/TF risks.
			\begin{itemize}
				\item CDD, EDD, record-keeping, internal controls.
				\item \textbf{Key:} Measures must be \textit{proportionate} to the identified risk.
			\end{itemize}
		\end{definition}
		\begin{alertblock}{Exam Focus}
			\textbf{IO.3 and IO.4 are prioritized} by supervisors evaluating bank compliance. Know which IOs correspond to national risk assessments, supervision, and preventive measures.
		\end{alertblock}
	\end{frame}
	
	% ============================================================
	% SLIDE 5: DORA - REGULATORY EXPECTATIONS
	% ============================================================
	\begin{frame}{Digital Operational Resilience Act (DORA)}
		\begin{block}{Scope}
			Applies to all financial entities operating in the EU, including banks, investment firms, crypto-asset service providers, and third-party ICT providers.
		\end{block}
		\begin{block}{Key Regulatory Expectation}
			\begin{center}
				\textbf{Centralizing oversight} of Information and Communications Technology (ICT) risk management across all member-state branches.
			\end{center}
			\begin{itemize}
				\item Financial institutions must have a \textbf{unified ICT risk governance framework}.
				\item Conduct \textbf{ICT risk assessments} and implement \textbf{incident reporting mechanisms}.
				\item Manage \textbf{third-party ICT provider risks} through contracts and continuous monitoring.
			\end{itemize}
		\end{block}
		\begin{alertblock}{Critical Action for Cross-Border Banks}
			Relying solely on national regulators' cybersecurity requirements is \textbf{insufficient}. DORA mandates centralized, pan-European oversight.
		\end{alertblock}
	\end{frame}
	
	% ============================================================
	% SLIDE 6: EARLY EU AML DIRECTIVES - SHORTCOMINGS
	% ============================================================
	\begin{frame}{Early EU AML Directives (1st \& 2nd)}
		\begin{block}{Major Shortcoming}
			\begin{center}
				\textbf{Failure to ensure timely and complete transposition} by all member states.
			\end{center}
		\end{block}
		\begin{itemize}
			\item \textbf{1st Directive (1991):} Focused on drug trafficking; lacked uniformity in implementation.
			\item \textbf{2nd Directive (2001):} Expanded predicate offenses but still suffered from divergent national interpretations.
			\item \textbf{Consequence:} Regulatory arbitrage—criminals exploited weak AML frameworks in certain EU countries.
		\end{itemize}
		\begin{block}{Later Revisions (3rd-6th Directives)}
			\begin{itemize}
				\item \textbf{3rd (2005):} Incorporated FATF Recommendations; introduced beneficial ownership requirements.
				\item \textbf{4th (2015):} Required central registers of beneficial owners.
				\item \textbf{5th (2018):} Enhanced transparency; extended to virtual currencies and prepaid cards.
				\item \textbf{6th (2020):} Harmonized predicate offenses and criminal penalties across member states.
			\end{itemize}
		\end{block}
	\end{frame}
	
	% ============================================================
	% SLIDE 7: INTERNAL AFC REPORTS - GOVERNANCE CONTENT
	% ============================================================
	\begin{frame}{Internal AFC Reporting to Governance Committees}
		\begin{block}{Appropriate Contents (Select Three)}
			\begin{enumerate}
				\item \textbf{Updates on regulatory interactions and supervisory reviews.}
				\begin{itemize}
					\item Include findings from on-site inspections, regulatory letters, and enforcement actions.
				\end{itemize}
				\item \textbf{Assessment outcomes of control effectiveness on monitoring tools.}
				\begin{itemize}
					\item Report on key performance indicators (KPIs) and key risk indicators (KRIs).
					\item Identify gaps in transaction monitoring, screening, and reporting systems.
				\end{itemize}
				\item \textbf{Investigation statuses of escalated SAR reviews.}
				\begin{itemize}
					\item Provide metrics on SAR filings, case disposition, and law enforcement feedback.
				\end{itemize}
			\end{enumerate}
		\end{block}
		\begin{alertblock}{What NOT to Include}
			\textbf{Cash flow performance figures segmented by region} are \textit{not} a core AFC governance report component—these are financial performance metrics, not AML risk indicators.
		\end{alertblock}
	\end{frame}
	
	% ============================================================
	% SLIDE 8: AML POLICY - MANDATORY ELEMENTS
	% ============================================================
	\begin{frame}{AML/CFT Policy Mandatory Elements}
		\begin{block}{Key Elements (Select Two)}
			\begin{enumerate}
				\item \textbf{Assignment of compliance responsibilities.}
				\begin{itemize}
					\item Must designate a \textbf{Compliance Officer (MLRO in the UK)} with clear authority and resources.
					\item Define roles and responsibilities across the three lines of defense.
				\end{itemize}
				\item \textbf{Program review triggers} based on regulatory or business events.
				\begin{itemize}
					\item Reviews must be triggered by: new regulations, significant business changes, new products/services, or changes in risk appetite.
					\item Frequency: Typically at least annually, plus event-driven reviews.
				\end{itemize}
			\end{enumerate}
		\end{block}
		\begin{block}{Other Essential Components}
			\begin{itemize}
				\item \textbf{Risk assessment framework} (enterprise-wide).
				\item \textbf{Internal controls and audit} requirements.
				\item \textbf{Record-keeping and reporting} obligations.
			\end{itemize}
		\end{block}
	\end{frame}
	
	% ============================================================
	% SLIDE 9: ENTERPRISE-WIDE RISK ASSESSMENT (EWRA)
	% ============================================================
	\begin{frame}{Enterprise-Wide Risk Assessment (EWRA)}
		\begin{block}{What is EWRA?}
			A comprehensive, documented risk assessment that identifies, assesses, and mitigates ML/TF risks across the entire organization—not just individual business units.
		\end{block}
		\begin{block}{Scenario: High-Risk Jurisdiction Customers via Mobile}
			The institution must \textbf{incorporate digital delivery channels} into the EWRA to reassess risk exposure.
			\begin{itemize}
				\item Mobile platforms often bypass branch monitoring, creating a \textbf{risk blind spot}.
				\item Must assess the \textbf{risk of customers from high-risk jurisdictions} using digital channels.
				\item \textbf{Action:} Update the EWRA to include digital-specific risk indicators.
			\end{itemize}
		\end{block}
		\begin{alertblock}{Regulatory Expectation}
			EWRA must be \textbf{dynamic}—updated when new channels, products, or customer segments emerge.
		\end{alertblock}
	\end{frame}
	
	% ============================================================
	% SLIDE 10: TCSP - HIGH-RISK CLIENT HANDLING
	% ============================================================
	\begin{frame}{Trust and Company Service Providers (TCSPs)}
		\begin{block}{Scenario}
			A TCSP firm is reviewing whether to accept a client who wants to incorporate in a low-tax Caribbean jurisdiction. The potential client requests confidentiality protections for all shareholders. A connection is identified to industries marked high-risk in the NRA.
		\end{block}
		\begin{block}{Correct Response}
			\begin{center}
				\textbf{Decline immediately} unless full beneficial owner information and risk justifications are disclosed and verified.
			\end{center}
		\end{block}
		\begin{itemize}
			\item Connections to high-risk industries (from the National Risk Assessment) and confidentiality requests are \textbf{significant red flags}.
			\item Offloading responsibility to "local counsel" is \textbf{not} a valid risk mitigation strategy.
			\item \textbf{Alternative:} Refer to internal high-risk committee for EDD before acceptance.
		\end{itemize}
	\end{frame}
	
	% ============================================================
	% SLIDE 11: PEPs - ENHANCED SCRUTINY (SINGAPORE)
	% ============================================================
	\begin{frame}{Politically Exposed Persons (PEPs) - Singapore Framework}
		\begin{block}{Scenario}
			A compliance officer at a financial institution in Singapore is reviewing customer onboarding files. The customer is a PEP from outside the region, with a complex ownership structure and high-value international transactions.
		\end{block}
		\begin{block}{Correct Action}
			\begin{center}
				Refer the case to the \textbf{high-risk customer review committee} for enhanced due diligence (EDD) before onboarding.
			\end{center}
		\end{block}
		\begin{itemize}
			\item Singapore MAS requires EDD for all PEPs, especially those from outside the region.
			\item EDD includes: source of wealth, source of funds, and ongoing monitoring.
			\item \textbf{Complex ownership structures} require beneficial ownership tracing.
			\item \textbf{High-value international transactions} trigger additional scrutiny.
		\end{itemize}
	\end{frame}
	
	% ============================================================
	% SLIDE 12: MILITARY ORGANIZATIONS - RISK FACTORS
	% ============================================================
	\begin{frame}{Military Organizations and Financial Crime Risks}
		\begin{block}{Risk Elements (Select Two)}
			\begin{enumerate}
				\item \textbf{Government ownership of defense contractors.}
				\begin{itemize}
					\item Creates risk of \textbf{state capture, corruption, and diversion of public funds.}
					\item Defense procurement is often non-transparent.
				\end{itemize}
				\item \textbf{Oversight by PEPs.}
				\begin{itemize}
					\item Military officials often qualify as PEPs (especially foreign military officers).
					\item Risk of bribery, illicit payments, and sanctions evasion.
				\end{itemize}
			\end{enumerate}
		\end{block}
		\begin{alertblock}{Note on Distractor}
			"Clear regulations for export of dual-use technology" is \textit{not} a risk factor—it's a control measure. Similarly, "oversight by PEPs" is correctly identified as a risk.
		\end{alertblock}
	\end{frame}
	
	% ============================================================
	% SLIDE 13: DATA COVERAGE GAP ASSESSMENT
	% ============================================================
	\begin{frame}{Data Coverage and Gap Assessment Challenges}
		\begin{block}{What is a Data Gap Assessment?}
			A process to identify deficiencies in AML data sources, quality, and integration.
		\end{block}
		\begin{block}{Challenges Identified (Select Two)}
			\begin{enumerate}
				\item \textbf{Missing identifiers in customer profiles.}
				\begin{itemize}
					\item E.g., missing DOB, passport numbers, or beneficial ownership information.
					\item Impacts the ability to screen against sanctions and PEP lists.
				\end{itemize}
				\item \textbf{Inconsistent use of country codes across systems.}
				\begin{itemize}
					\item Example: "US" vs. "USA" vs. "840" vs. "United States."
					\item Leads to false negatives and false positives in monitoring and reporting.
				\end{itemize}
			\end{enumerate}
		\end{block}
		\begin{alertblock}{Consequences}
			Data gaps and inconsistencies undermine the effectiveness of transaction monitoring, screening, and regulatory reporting.
		\end{alertblock}
	\end{frame}
	
	% ============================================================
	% SLIDE 14: DATA CLEANSING - DEFINITION AND PROCESS
	% ============================================================
	\begin{frame}{Data Cleansing Processes Explained}
		\begin{block}{Definition}
			Data cleansing (or data cleaning) is the process of identifying and correcting (or removing) inaccurate, incomplete, or duplicated data.
		\end{block}
		\begin{block}{Examples (Select Two)}
			\begin{enumerate}
				\item \textbf{Removal of duplicate entries.}
				\begin{itemize}
					\item Duplicate profiles in merged databases create risk of incomplete due diligence.
					\item Example: The same customer appears twice with different addresses and transaction histories.
				\end{itemize}
				\item \textbf{Validation of data accuracy.}
				\begin{itemize}
					\item Checking that data conforms to defined formats (e.g., date formats, phone numbers, addresses).
					\item Example: Validating that a "ZIP code" field contains valid postal codes.
				\end{itemize}
			\end{enumerate}
		\end{block}
		\begin{block}{Why It Matters}
			Data cleansing is a \textbf{pre-requisite} for effective AML controls. Without clean data, transaction monitoring alerts are unreliable.
		\end{block}
	\end{frame}
	
	% ============================================================
	% SLIDE 15: DATA VALIDATION AND TESTING - AFC SYSTEMS
	% ============================================================
	\begin{frame}{Data Validation and Testing in AFC Systems}
		\begin{block}{Objectives (Select Three)}
			\begin{enumerate}
				\item \textbf{Ensuring data inputs reflect actual customer and transaction behaviors.}
				\begin{itemize}
					\item Test that data from source systems (core banking, CRM, etc.) accurately feeds into AML systems.
				\end{itemize}
				\item \textbf{Verifying that automated controls are functioning correctly based on input quality.}
				\begin{itemize}
					\item Validate that transaction monitoring rules fire alerts correctly.
					\item Test that screening systems block or alert on prohibited names.
				\end{itemize}
				\item \textbf{Testing whether controls generate alerts under appropriate risk conditions.}
				\begin{itemize}
					\item Use known test cases (e.g., hypothetical suspicious scenarios) to validate system logic.
					\item Ensure alerts are generated for true positives, not just false positives.
				\end{itemize}
			\end{enumerate}
		\end{block}
	\end{frame}
	
	% ============================================================
	% SLIDE 16: AI TOOLS - IMPLEMENTATION BEST PRACTICES
	% ============================================================
	\begin{frame}{Transitioning to AI-Based Compliance Tools}
		\begin{block}{Maintain Effectiveness (Select Three)}
			\begin{enumerate}
				\item \textbf{Ensuring data quality for training algorithms.}
				\begin{itemize}
					\item "Garbage in, garbage out"—AI models are only as good as their training data.
					\item Need to cleanse and normalize data before AI deployment.
				\end{itemize}
				\item \textbf{Managing the change in employee workflows.}
				\begin{itemize}
					\item AI will change how analysts work; must provide adequate training.
					\item Human-in-the-loop is still required for validation of AI outputs.
				\end{itemize}
				\item \textbf{Verifying transparency and explainability in AI models.}
				\begin{itemize}
					\item Regulators (e.g., EU GDPR, US Fair Lending) require \textit{explainable AI}.
					\item Must be able to justify why an AI model flagged a particular transaction or customer.
				\end{itemize}
			\end{enumerate}
		\end{block}
		\begin{alertblock}{Distractor}
			"Selecting third-party global audit providers" is not a primary concern for AI effectiveness—focus is on data, people, and explainability.
		\end{alertblock}
	\end{frame}
	
	% ============================================================
	% SLIDE 17: ML MODEL VALIDATION - REGULATORY REQUIREMENT
	% ============================================================
	\begin{frame}{Machine Learning Model Validation}
		\begin{block}{Scenario}
			An international bank is implementing a new machine learning tool to detect suspicious transaction patterns. The regulator in one of its jurisdictions requires validation and explainability of all models used in transaction monitoring.
		\end{block}
		\begin{block}{Correct Approach}
			\begin{center}
				\textbf{Document model development}, ensure explainability, and perform validation \textbf{before deployment}.
			\end{center}
			\begin{itemize}
				\item \textbf{Validation:} Test model on historical data to measure accuracy, precision, recall, and false positive rate.
				\item \textbf{Explainability:} Document model parameters, variables, and decision logic.
				\item \textbf{Documentation:} Maintain a comprehensive model risk management (MRM) file.
			\end{itemize}
		\end{block}
		\begin{alertblock}{Key Regulatory Concern}
			"Black box" models are not acceptable; regulators require \textit{understandable} and \textit{auditable} models.
		\end{alertblock}
	\end{frame}
	
	% ============================================================
	% SLIDE 18: ENTITY RESOLUTION IN BLOCKCHAIN ANALYTICS
	% ============================================================
	\begin{frame}{Entity Resolution in Blockchain Analytics}
		\begin{definition}[Entity Resolution]
			The process of \textbf{linking data from multiple sources} to improve match accuracy.
		\end{definition}
		\begin{block}{How It Works}
			\begin{itemize}
				\item Blockchain analytics tools (e.g., Chainalysis, Elliptic) use entity resolution to link:
				\begin{itemize}
					\item On-chain wallet addresses.
					\item Off-chain KYC data.
					\item Sanctions lists, PEP lists, and adverse media.
				\end{itemize}
				\item Example: A wallet address associated with a ransomware attack is linked to an exchange account through entity resolution.
			\end{itemize}
		\end{block}
		\begin{alertblock}{Exam Focus}
			Entity resolution is \textbf{not} just about matching names; it's about linking \textit{multiple identifiers} across systems to build a complete risk profile.
		\end{alertblock}
	\end{frame}
	
	% ============================================================
	% SLIDE 19: FUZZY LOGIC IN PEP SCREENING
	% ============================================================
	\begin{frame}{Fuzzy Logic in PEP Screening}
		\begin{block}{What is Fuzzy Logic?}
			A computational technique that accommodates \textbf{approximate, partial, or incomplete matches}—unlike exact matching (Boolean logic).
		\end{block}
		\begin{block}{Why Integrate Fuzzy Logic?}
			\begin{center}
				Because it accommodates \textbf{inconsistencies in data formats and naming conventions}.
			\end{center}
		\end{block}
		\begin{itemize}
			\item PEP lists often contain variations: "William" vs. "Bill," "Mohammed" vs. "Muhammad," "Smith" vs. "Smyth."
			\item Foreign names transliterated differently (e.g., Arabic or Chinese names).
			\item Fuzzy logic calculates a \textbf{similarity score} and flags matches above a certain threshold.
			\item \textbf{Challenge:} Setting the threshold correctly to avoid false positives and negatives.
		\end{itemize}
	\end{frame}
	
	% ============================================================
	% SLIDE 20: PASSIVE LIVENESS DETECTION - KYC
	% ============================================================
	\begin{frame}{Passive vs. Active Liveness Detection}
		\begin{block}{Passive Liveness Detection}
			\begin{center}
				Uses \textbf{minimal user interaction} and analyzes subtle facial movement or texture patterns.
			\end{center}
			\begin{itemize}
				\item Detects liveness by analyzing:
				\begin{itemize}
					\item Micro-expressions (e.g., blinking, slight head movements).
					\item Texture analysis (e.g., skin texture, depth perception).
					\item Spoof detection (e.g., detects if using a photo or video).
				\end{itemize}
				\item \textbf{Advantage:} Seamless user experience—no action required from the user.
			\end{itemize}
		\end{block}
		\begin{block}{Active Liveness Detection}
			\begin{itemize}
				\item Requires explicit user actions: blinking, smiling, turning head.
				\item More secure but more intrusive; may have lower completion rates.
			\end{itemize}
		\end{block}
		\begin{alertblock}{Exam Distractor}
			"Asks users to answer KYC security questions during video evaluation" is \textit{not} liveness detection—it's KYC verification.
		\end{alertblock}
	\end{frame}
	
	% ============================================================
	% SLIDE 21: ROUND-TRIP TRANSACTIONS - ANALYST ACTION
	% ============================================================
	\begin{frame}{Rules-Based Alerts: Round-Trip Transactions}
		\begin{block}{What is a "Round-Trip" Transaction?}
			A series of transactions where funds are transferred out and then returned to the same account, often within a short timeframe. Typically associated with \textbf{layering}.
		\end{block}
		\begin{block}{Analyst Workflow}
			\begin{enumerate}
				\item Review the alert—identify a nearly identical send and receive remittance within minutes.
				\item Confirm the transaction pattern and recent account activity.
				\item \textbf{Escalate to Level 2} for deeper investigation.
			\end{enumerate}
		\end{block}
		\begin{alertblock}{Why Escalate?}
			Level 1 analysts are trained for \textbf{initial triage and data gathering}. Level 2 analysts have the experience and authority to conduct full investigations, interview customers, and potentially file SARs.
		\end{alertblock}
	\end{frame}
	
	% ============================================================
	% SLIDE 22: CARD-BASED PRODUCT RISK FACTORS
	% ============================================================
	\begin{frame}{Card-Based Products: Risk Factors}
		\begin{block}{Common Risk Factors (Select Three)}
			\begin{enumerate}
				\item \textbf{Anonymity through prepaid cards.}
				\begin{itemize}
					\item Prepaid cards can be purchased with cash and used without identification.
					\item Reloadable cards can be used repeatedly without ongoing CDD.
				\end{itemize}
				\item \textbf{High volumes of rapid transactions.}
				\begin{itemize}
					\item Card transactions are processed instantly in high volumes.
					\item Structuring (smurfing) can occur using multiple cards across different merchants.
				\end{itemize}
				\item \textbf{Use by customers in high-risk sectors.}
				\begin{itemize}
					\item E.g., merchants in crypto, gambling, or cash-intensive businesses.
					\item These sectors generate high-risk card transactions requiring enhanced scrutiny.
				\end{itemize}
			\end{enumerate}
		\end{block}
	\end{frame}
	
	% ============================================================
	% SLIDE 23: PREPAID CARDS - VULNERABILITIES
	% ============================================================
	\begin{frame}{Prepaid Card Vulnerabilities}
		\begin{block}{Key Vulnerability}
			\begin{center}
				\textbf{Can be used anonymously, especially when reloadable or unregistered.}
			\end{center}
		\end{block}
		\begin{itemize}
			\item \textbf{Unregistered cards:} Can be purchased with cash at retail outlets; no customer identification required.
			\item \textbf{Reloadable cards:} Can be topped up multiple times, potentially with illicit funds.
			\item \textbf{Cross-border use:} Can be used internationally, bypassing local AML controls.
			\item \textbf{Regulatory Response:} Many jurisdictions (e.g., EU 5th AML Directive) have imposed limits on anonymous prepaid cards (e.g., €150 cap).
		\end{itemize}
		\begin{alertblock}{Exam Point}
			Know the specific vulnerabilities of \textit{prepaid}, \textit{reloadable}, and \textit{anonymous} cards. Each has different risk profiles.
		\end{alertblock}
	\end{frame}
	
	% ============================================================
	% SLIDE 24: E-COMMERCE TRANSACTION MONITORING
	% ============================================================
	\begin{frame}{E-Commerce AML Challenges}
		\begin{block}{The Challenge}
			\begin{center}
				\textbf{High transaction volumes and decentralized payments} make patterns harder to detect.
			\end{center}
		\end{block}
		\begin{itemize}
			\item E-commerce platforms process millions of small-value transactions daily.
			\item Payments are decentralized—multiple payment gateways, cryptocurrencies, and digital wallets.
			\item \textbf{Consequence:} Traditional rules-based monitoring generates too many false positives.
			\item \textbf{Mitigation:} Use AI/ML models to detect anomalies; segment data by merchant, payment method, and product.
		\end{itemize}
		\begin{block}{Additional Risks}
			\begin{itemize}
				\item \textbf{Goods/services as money laundering vehicles:} High-value goods (e.g., luxury items) purchased with illicit funds.
				\item \textbf{Trade-based money laundering (TBML):} Mis-invoicing, over/under-invoicing.
			\end{itemize}
		\end{block}
	\end{frame}
	
	% ============================================================
	% SLIDE 25: HIGH-VALUE ASSETS AND CRYPTO
	% ============================================================
	\begin{frame}{High-Value Assets and Cryptocurrency}
		\begin{block}{Red Flag}
			\begin{center}
				Making payments for high-value items in \textbf{cryptocurrency}.
			\end{center}
		\end{block}
		\begin{itemize}
			\item Cryptocurrency offers \textbf{pseudonymity}, making it attractive for illicit payments.
			\item High-value items (e.g., real estate, art, luxury cars) purchased with crypto raise suspicion.
			\item \textbf{Key Challenge:} Volatility and price fluctuations complicate source-of-funds verification.
		\end{itemize}
		\begin{block}{Additional Red Flags for Real Estate}
			\begin{itemize}
				\item \textbf{Foreign PEPs} purchasing multiple properties via trust companies.
				\item \textbf{Complex financing plans} spanning several years may be a red flag, but \textit{cryptocurrency payments} are more directly associated with laundering.
				\item \textbf{Trusts and shell companies} used to obscure beneficial ownership.
			\end{itemize}
		\end{block}
	\end{frame}
	
	% ============================================================
	% SLIDE 26: FREE-TRADE ZONES (FTZs) - VULNERABILITIES
	% ============================================================
	\begin{frame}{Free-Trade Zone Vulnerabilities}
		\begin{block}{Typical Characteristics (Select Two)}
			\begin{enumerate}
				\item \textbf{Limited customs authority} over transitory goods.
				\begin{itemize}
					\item Goods can move in and out with minimal documentation.
					\item Creates opportunity for mis-invoicing, smuggling, and trade-based ML.
				\end{itemize}
				\item \textbf{Evasions of customs payments.}
				\begin{itemize}
					\item Companies may declare goods for re-export to avoid duties, then divert them to the local market.
					\item Example: Imports electronics marked for re-export; actually sells them locally, evading customs duties.
				\end{itemize}
			\end{enumerate}
		\end{block}
		\begin{alertblock}{Exam Distractor}
			"Laundering through over-invoicing" is a TBML technique, but the specific FTZ risk in the scenario (electronics for re-export) is \textit{evasion of customs payments}.
		\end{alertblock}
	\end{frame}
	
	% ============================================================
	% SLIDE 27: PRIVATE BANKING - CORE VULNERABILITIES
	% ============================================================
	\begin{frame}{Private Banking \& Wealth Management}
		\begin{block}{Core ML Vulnerability}
			\begin{center}
				\textbf{Discretion and high client autonomy} can obscure beneficial ownership and reduce transparency.
			\end{center}
		\end{block}
		\begin{itemize}
			\item Private banking clients expect \textbf{confidentiality}—this can be abused to hide illicit activities.
			\item \textbf{Complex structures:} Trusts, foundations, offshore companies, and nominee arrangements.
			\item \textbf{High autonomy:} Clients may execute transactions directly, bypassing internal controls.
			\item \textbf{Relationship managers:} May become overly close to clients, failing to report suspicious activity.
		\end{itemize}
		\begin{block}{Mitigations}
			\begin{itemize}
				\item Enhanced due diligence (EDD) for high-net-worth individuals and PEPs.
				\item Independent review of client transactions.
				\item Regular training for relationship managers on red flags.
			\end{itemize}
		\end{block}
	\end{frame}
	
	% ============================================================
	% SLIDE 28: INVESTMENT BANKING - UNUSUAL ACTIVITY
	% ============================================================
	\begin{frame}{Investment Banking: Sudden Business Changes}
		\begin{block}{Red Flag Scenario}
			A corporate investment bank client typically deals in low-margin bond offerings but suddenly begins issuing high-return instruments through an unfamiliar subsidiary.
		\end{block}
		\begin{block}{Correct Response}
			\begin{center}
				\textbf{Escalate internally} and request further documentation before determining next steps.
			\end{center}
		\end{block}
		\begin{itemize}
			\item Sudden departure from normal business profile is a \textbf{red flag}.
			\item Unfamiliar subsidiary may be a shell or front company.
			\item \textbf{Action:} Obtain additional due diligence on the subsidiary, including beneficial ownership, financial statements, and rationale for the new product.
		\end{itemize}
	\end{frame}
	
	% ============================================================
	% SLIDE 29: BANK SECRECY ACT (BSA) - PURPOSE
	% ============================================================
	\begin{frame}{Bank Secrecy Act (BSA) / FinCEN}
		\begin{block}{Purpose of FinCEN}
			\begin{center}
				Safeguards the financial system by administering and enforcing the Bank Secrecy Act to \textbf{detect and prevent financial crimes}.
			\end{center}
		\end{block}
		\begin{itemize}
			\item FinCEN is the U.S. Financial Intelligence Unit (FIU).
			\item Administers BSA regulations and collects Suspicious Activity Reports (SARs).
			\item Issues guidance and imposes civil penalties for non-compliance.
			\item \textbf{Key BSA Requirements:}
			\begin{itemize}
				\item CTRs for cash transactions over \$10,000.
				\item SARs for suspicious activity over \$5,000 (or \$2,000 for MSBs).
				\item Record-keeping requirements for funds transfers.
			\end{itemize}
		\end{itemize}
	\end{frame}
	
	% ============================================================
	% SLIDE 30: CONSUMER PROTECTION - AML INTERSECTION
	% ============================================================
	\begin{frame}{Consumer Protection in Financial Services}
		\begin{block}{Primary Objectives (Select Two)}
			\begin{enumerate}
				\item \textbf{Prevent unauthorized transactions} affecting vulnerable customers.
				\begin{itemize}
					\item Consumer protection and AML overlap—fraud prevention is a core AML obligation.
					\item Vulnerable customers include elderly, minors, and individuals with limited financial literacy.
				\end{itemize}
				\item \textbf{Ensure clear and transparent disclosures} of terms.
				\begin{itemize}
					\item Customers must understand fees, risks, and obligations.
					\item Transparency reduces opportunities for mis-selling and financial abuse.
				\end{itemize}
			\end{enumerate}
		\end{block}
		\begin{alertblock}{Distractor}
			"Promote transaction surveillance to reduce bank overhead" is \textit{not} a consumer protection objective—it's an efficiency goal.
		\end{alertblock}
	\end{frame}
	
	% ============================================================
	% SLIDE 31: REGULATORY REPORTING INCONSISTENCIES
	% ============================================================
	\begin{frame}{AFC Regulatory Report Inconsistencies}
		\begin{block}{Scenario}
			A financial institution regularly receives regulatory feedback on its AFC regulatory report presentation being inconsistent across filings. One quarter, the report uses free-text formatting; the next quarter, it uses structured tabular format.
		\end{block}
		\begin{block}{Impact}
			\begin{center}
				The institution \textbf{reduces the efficiency of regulatory review}.
			\end{center}
		\end{block}
		\begin{itemize}
			\item \textbf{Why it matters:} Regulators rely on consistent, machine-readable formats to review data efficiently.
			\item Mixed formats force regulators to re-map data each time, increasing review time and error risk.
			\item \textbf{Solution:} Standardize reporting templates across all business units and jurisdictions.
		\end{itemize}
	\end{frame}
	
	% ============================================================
	% SLIDE 32: DATA PRIVACY COMPLIANCE - CROSS-BORDER
	% ============================================================
	\begin{frame}{Data Privacy Compliance Across Jurisdictions}
		\begin{block}{Scenario}
			A global compliance team is updating its investigation protocols to conform to data privacy laws across jurisdictions (e.g., GDPR, CCPA, PDPA).
		\end{block}
		\begin{block}{Correct Actions (Select Two)}
			\begin{enumerate}
				\item \textbf{Create jurisdiction-specific procedures} for capturing research findings, including required formats and approval controls.
				\begin{itemize}
					\item GDPR requires data minimization, purpose limitation, and lawful processing.
					\item Each jurisdiction may have unique approval requirements for cross-border data transfers.
				\end{itemize}
				\item \textbf{Develop a standard research template with adjustable sections} based on local data regulations.
				\begin{itemize}
					\item Core sections are consistent (e.g., investigation findings, rationale).
					\item Adjustable sections accommodate local legal requirements (e.g., retention periods, consent forms).
				\end{itemize}
			\end{enumerate}
		\end{block}
	\end{frame}
	
	% ============================================================
	% SLIDE 33: MSB CLASSIFICATION - COMPLIANCE OBLIGATIONS
	% ============================================================
	\begin{frame}{Money Services Businesses (MSBs)}
		\begin{block}{Why Proper Classification Matters}
			\begin{center}
				MSB designation triggers \textbf{specific AML compliance and reporting obligations}.
			\end{center}
		\end{block}
		\begin{itemize}
			\item In the U.S., FinCEN regulates MSBs under the BSA.
			\item Obligations include:
			\begin{itemize}
				\item Register with FinCEN.
				\item Implement a written AML program.
				\item File CTRs and SARs.
				\item Maintain records of funds transfers and MSB registration.
			\end{itemize}
			\item \textbf{Consequence of misclassification:} Failure to meet MSB obligations results in fines and enforcement actions.
		\end{itemize}
	\end{frame}
	
	% ============================================================
	% SLIDE 34: PSPs - CROSS-BORDER STRATEGY
	% ============================================================
	\begin{frame}{Payment Service Providers (PSPs)}
		\begin{block}{Scenario}
			A PSP is planning to offer new services targeting small merchants in multiple countries, some of which lack mature AML frameworks.
		\end{block}
		\begin{block}{Correct Strategy}
			\begin{center}
				\textbf{Customize AML policy and onboarding procedures} per jurisdiction and product.
			\end{center}
		\end{block}
		\begin{itemize}
			\item One-size-fits-all approach fails to address jurisdictional risk differences.
			\item \textbf{High-risk jurisdictions:} Enhanced due diligence, transaction limits, or restricted services.
			\item \textbf{Low-risk jurisdictions:} Simplified due diligence (SDD) may be permitted.
			\item \textbf{Product-based customization:} E.g., low-value vs. high-value products have different risk profiles.
		\end{itemize}
		\begin{alertblock}{Distractor}
			"Apply home-country AML standards across all jurisdictions" is \textit{not} optimal—it may be overly prescriptive or insufficient in different contexts.
		\end{alertblock}
	\end{frame}
	
	% ============================================================
	% SLIDE 35: CASINO AML CONTROLS
	% ============================================================
	\begin{frame}{Casino AML Controls}
		\begin{block}{Control Measures (Select Two)}
			\begin{enumerate}
				\item \textbf{Enforcing chip redemption policies} tied to minimum gaming thresholds.
				\begin{itemize}
					\item Prevents "chip washing"—exchanging chips for cash without meaningful gaming activity.
					\item Common threshold: \$10,000 in the U.S. triggers reporting.
				\end{itemize}
				\item \textbf{Tracking player activity} through registered loyalty programs.
				\begin{itemize}
					\item Loyalty programs provide rich data on player behavior.
					\item Enables tracking of wins, losses, and session durations to detect anomalies.
				\end{itemize}
			\end{enumerate}
		\end{block}
		\begin{alertblock}{Distractor}
			"Monitoring unstructured table gameplay for win/loss patterns" is too granular and not a formal control measure—it's a \textit{monitoring technique}.
		\end{alertblock}
	\end{frame}
	
	% ============================================================
	% SLIDE 36: CANNABIS BUSINESSES - LEGAL COMPLEXITY
	% ============================================================
	\begin{frame}{Cannabis-Related Businesses (CRBs)}
		\begin{block}{Scenario}
			A client wants to open a retail dispensary account in a province/state where recreational marijuana is legalized, even though marijuana remains illegal under national law. The financial institution is federally regulated.
		\end{block}
		\begin{block}{Correct Action}
			\begin{center}
				Apply \textbf{enhanced due diligence}, assess provincial/state law compliance, and consult internal legal guidance.
			\end{center}
		\end{block}
		\begin{itemize}
			\item \textbf{Federal vs. State/Provincial Law:} Federal law (e.g., U.S. Controlled Substances Act) prohibits marijuana; state law may permit it.
			\item \textbf{FinCEN Guidance:} Requires filing SARs on CRB activity (Marijuana Limited SAR or Priority SAR).
			\item \textbf{EDD Elements:} Verify state licenses, monitor for cross-state activities, and assess source of funds.
		\end{itemize}
	\end{frame}
	
	% ============================================================
	% SLIDE 37: CRYPTO SANCTIONS SCREENING
	% ============================================================
	\begin{frame}{Cryptocurrency and Sanctions Compliance}
		\begin{block}{Minimizing Sanctions Risk}
			\begin{center}
				Implement robust \textbf{blockchain analytics} and transaction monitoring tools to screen wallet addresses against sanctions lists.
			\end{center}
		\end{block}
		\begin{itemize}
			\item \textbf{Sanctions screening:} Must screen all crypto transactions against OFAC, EU, UN, and other sanctions lists.
			\item \textbf{Entity resolution:} Links on-chain addresses to off-chain risk data (e.g., identifies wallets associated with sanctioned entities).
			\item \textbf{Real-time screening:} Transactions must be blocked or rejected before settlement.
			\item \textbf{Travel Rule compliance:} FATF Recommendation 16 requires transfer of originator/beneficiary info for crypto transfers.
		\end{itemize}
	\end{frame}
	
	% ============================================================
	% SLIDE 38: PERPETUAL KYC - TECHNOLOGY SOLUTION
	% ============================================================
	\begin{frame}{Perpetual KYC (pKYC)}
		\begin{block}{Technology Solution}
			\begin{center}
				Implement \textbf{perpetual KYC} with centralized data and trigger-based updates.
			\end{center}
		\end{block}
		\begin{itemize}
			\item \textbf{What is pKYC?} Continuous, event-driven KYC refresh rather than periodic reviews (e.g., every 3-5 years).
			\item \textbf{Triggers for pKYC Updates:}
			\begin{itemize}
				\item Changes in PEP status.
				\item Adverse media hits.
				\item Unusual transaction patterns.
				\item Jurisdictional risk changes.
			\end{itemize}
			\item \textbf{Benefits:} Reduces customer friction, improves accuracy, and ensures timely EDD.
		\end{itemize}
	\end{frame}
	
	% ============================================================
	% SLIDE 39: DATA GOVERNANCE AND QUALITY CHALLENGES
	% ============================================================
	\begin{frame}{Data Governance and AML Challenges}
		\begin{block}{Common Issues in AML Data}
			\begin{itemize}
				\item \textbf{Missing identifiers} in customer profiles (e.g., no DOB, no tax ID).
				\item \textbf{Inconsistent address formats} (e.g., "123 Main St." vs. "123 Main Street").
				\item \textbf{Duplicate profiles} resulting from system mergers or poor data entry.
				\item \textbf{Misalignment} between transaction monitoring scenarios and typologies (e.g., rules detecting the wrong behaviors).
			\end{itemize}
		\end{block}
		\begin{block}{Remedy}
			\begin{center}
				Initiate a \textbf{data cleansing process} targeting missing and duplicated records.
			\end{center}
		\end{block}
	\end{frame}
	
	% ============================================================
	% SLIDE 40: STAKEHOLDER MAPPING
	% ============================================================
	\begin{frame}{Stakeholder Management in Compliance}
		\begin{block}{Who are the Key Stakeholders?}
			\begin{itemize}
				\item \textbf{Internal:}
				\begin{itemize}
					\item Board of Directors and Audit Committee.
					\item Compliance Committee and MLRO.
					\item Business Units (e.g., Retail, Corporate, Wealth).
					\item IT and Data Management Teams.
					\item Internal Audit.
				\end{itemize}
				\item \textbf{External:}
				\begin{itemize}
					\item Regulators (e.g., FinCEN, FCA, MAS, ESMA).
					\item FIUs.
					\item Law Enforcement.
					\item Customers.
				\end{itemize}
			\end{itemize}
		\end{block}
		\begin{alertblock}{Key Principle}
			Effective stakeholder management requires \textbf{clear communication}, \textbf{escalation protocols}, and \textbf{regular reporting}.
		\end{alertblock}
	\end{frame}
	
	% ============================================================
	% SLIDE 41: COMPLIANCE CULTURE AND TRAINING
	% ============================================================
	\begin{frame}{Building a Strong Compliance Culture}
		\begin{block}{Elements of a Strong Culture}
			\begin{itemize}
				\item \textbf{Tone from the Top:} Leadership must actively prioritize and demonstrate commitment to AML/CFT.
				\item \textbf{Training:}
				\begin{itemize}
					\item Role-specific training tailored to each business unit.
					\item Case studies and scenario-based learning.
					\item Annual refresh and updates for regulatory changes.
				\end{itemize}
				\item \textbf{Awareness:} Employees should understand the "why" behind controls—not just the "what."
				\item \textbf{Whistleblower policies:} Encourage reporting without fear of retaliation.
			\end{itemize}
		\end{block}
	\end{frame}
	
	% ============================================================
	% SLIDE 42: POST-ACQUISITION AML INTEGRATION
	% ============================================================
	\begin{frame}{Post-Acquisition Integration - AML Strategy}
		\begin{block}{Scenario}
			During post-acquisition integration, a compliance team is reviewing merged customer databases and identifies multiple entries with inconsistent address formats and duplicate profiles.
		\end{block}
		\begin{block}{Correct Priority}
			\begin{center}
				Initiate a \textbf{data cleansing process} targeting missing and duplicated records.
			\end{center}
		\end{block}
		\begin{itemize}
			\item Immediate action is to \textbf{cleanse the data} to ensure accurate reporting, monitoring, and screening.
			\item Additional steps:
			\begin{enumerate}
				\item Conduct a \textbf{data coverage and gap assessment}.
				\item Harmonize \textbf{country code standards} and data formats.
				\item Align \textbf{transaction monitoring scenarios} across merged entities.
			\end{enumerate}
		\end{itemize}
	\end{frame}
	
	% ============================================================
	% SLIDE 43: FUTURE TRENDS - AFC
	% ============================================================
	\begin{frame}{Future Trends in Anti-Financial Crime}
		\begin{itemize}
			\item \textbf{AI and Machine Learning:}
			\begin{itemize}
				\item Enhanced detection of complex patterns (e.g., network analysis for syndicates).
				\item Predictive models to identify emerging risks.
			\end{itemize}
			\item \textbf{Blockchain Analytics:}
			\begin{itemize}
				\item Deeper on-chain monitoring for DeFi, NFTs, and cross-chain transactions.
				\item Enhanced entity resolution for attribution.
			\end{itemize}
			\item \textbf{Data Sharing:}
			\begin{itemize}
				\item Public-private partnerships for information exchange (e.g., JMLIT, FinCEN Exchange).
				\item Balance between privacy and security.
			\end{itemize}
			\item \textbf{RegTech:}
			\begin{itemize}
				\item Automation of compliance, reporting, and KYC.
				\item AI-driven regulatory change management.
			\end{itemize}
		\end{itemize}
	\end{frame}
	
	% ============================================================
	% SLIDE 44: GREY LIST REMEDIATION STEPS
	% ============================================================
	\begin{frame}{Remediating FATF Grey List Status}
		\begin{block}{Key Steps for Jurisdictions}
			\begin{enumerate}
				\item \textbf{Address strategic deficiencies} identified by FATF (e.g., in AML/CFT framework, supervision, or enforcement).
				\item Improve supervision of \textbf{DNFBPs} (e.g., casinos, real estate agents, lawyers).
				\item Enhance \textbf{beneficial ownership transparency}.
				\item Demonstrate \textbf{effective implementation} and enforcement of AML/CFT measures.
				\item Engage in \textbf{continuous reporting} to FATF on progress.
			\end{enumerate}
		\end{block}
		\begin{alertblock}{Consequences of Non-Remediation}
			Enhanced due diligence requirements for transactions involving the jurisdiction, reputational risk, and potential sanctions.
		\end{alertblock}
	\end{frame}
	
	% ============================================================
	% SLIDE 45: KEY TAKEAWAYS - CAMS ESSENTIALS
	% ============================================================
	\begin{frame}{Key Takeaways for CAMS Exam}
		\begin{enumerate}
			\item \textbf{Know Your Customer (KYC)} is the foundation of AML—CDD, EDD, and SDD.
			\item \textbf{Risk-based approach} is key: assess, mitigate, monitor, and adapt.
			\item \textbf{Data quality} is non-negotiable for effective AML controls.
			\item \textbf{FATF standards} are the global benchmark—know the 40 Recommendations and Immediate Outcomes.
			\item \textbf{Technology is a tool, not a replacement} for human judgment and oversight.
			\item \textbf{Regulations are dynamic}—stay updated on DORA, EU Directives, and national laws.
		\end{enumerate}
	\end{frame}
	
	% ============================================================
	% SLIDE 46: EXAM STRATEGY
	% ============================================================
	\begin{frame}{ACAMS Exam Strategy}
		\begin{itemize}
			\item \textbf{Read questions carefully:} Look for qualifiers like "Select Three," "Select Two," or "Select All That Apply."
			\item \textbf{Think like a compliance officer:}
			\begin{itemize}
				\item When in doubt, \textit{escalate}, \textit{document}, and \textit{validate}.
				\item Focus on the risk-based approach—not prescriptive rules.
			\end{itemize}
			\item \textbf{Understand the "why":} Don't just memorize definitions; understand the underlying logic.
			\item \textbf{Know the frameworks:} FATF, DORA, Wolfsberg, EU Directives, and national laws (e.g., BSA, MAS).
			\item \textbf{Time management:} 2 minutes per question; don't get stuck.
		\end{itemize}
	\end{frame}
	
	% ============================================================
	% SLIDE 47: SAMPLE QUESTION 1 - CARD RISK FACTORS
	% ============================================================
	\begin{frame}{Sample Question: Card-Based Product Risks}
		\begin{block}{Question}
			Which of the following represent common risk factors associated with card-based products? (Select Three.)
		\end{block}
		\begin{block}{Answer}
			\begin{enumerate}
				\item Anonymity through prepaid cards.
				\item High volumes of rapid transactions.
				\item Use by customers in high-risk sectors.
			\end{enumerate}
		\end{block}
		\begin{alertblock}{Explanation}
			These factors directly enable money laundering and terrorist financing. The anonymity of prepaid cards, rapid transaction volumes, and connection to high-risk sectors all increase AML risk.
		\end{alertblock}
	\end{frame}
	
	% ============================================================
	% SLIDE 48: SAMPLE QUESTION 2 - PEP SINGAPORE
	% ============================================================
	\begin{frame}{Sample Question: PEP Onboarding (Singapore)}
		\begin{block}{Question}
			A compliance officer at a financial institution in Singapore is reviewing a PEP case. What action is most consistent with Singapore's AML framework?
		\end{block}
		\begin{block}{Answer}
			Refer the case to the high-risk customer review committee for enhanced due diligence before onboarding.
		\end{block}
		\begin{alertblock}{Explanation}
			Singapore MAS requires EDD for all PEPs, especially those from outside the region. Committee-level approval ensures appropriate oversight and documentation.
		\end{alertblock}
	\end{frame}
	
	% ============================================================
	% SLIDE 49: SAMPLE QUESTION 3 - DATA CLEANSING
	% ============================================================
	\begin{frame}{Sample Question: Data Cleansing}
		\begin{block}{Question}
			Which actions are examples of data cleansing processes? (Select Two.)
		\end{block}
		\begin{block}{Answer}
			\begin{enumerate}
				\item Removal of duplicate entries.
				\item Validation of data accuracy.
			\end{enumerate}
		\end{block}
		\begin{alertblock}{Explanation}
			Data cleansing is about \textit{correcting and standardizing} data—not about \textit{gathering} data. Duplicate removal and accuracy validation are core cleansing activities.
		\end{alertblock}
	\end{frame}
	
	% ============================================================
	% SLIDE 50: SAMPLE QUESTION 4 - E-COMMERCE
	% ============================================================
	\begin{frame}{Sample Question: E-Commerce Monitoring}
		\begin{block}{Question}
			Why is transaction monitoring particularly challenging in an e-commerce risk context?
		\end{block}
		\begin{block}{Answer}
			High transaction volumes and decentralized payments make patterns harder to detect.
		\end{block}
		\begin{alertblock}{Explanation}
			Traditional rules-based monitoring struggles with high-volume, low-value transactions. Decentralized payments (multiple gateways, crypto) further complicate detection.
		\end{alertblock}
	\end{frame}
	
	% ============================================================
	% SLIDE 51: SAMPLE QUESTION 5 - FTZ RISK
	% ============================================================
	\begin{frame}{Sample Question: Free-Trade Zone Risks}
		\begin{block}{Question}
			A firm in a FTZ imports electronics marked for re-exportation. The firm also deals in the same electronics on the local market. What risk does this represent?
		\end{block}
		\begin{block}{Answer}
			\textbf{Evasion of customs payments.}
		\end{block}
		\begin{alertblock}{Explanation}
			The firm is likely declaring goods for re-export to avoid import duties, then diverting them to the local market—this is customs duty evasion, not over-invoicing.
		\end{alertblock}
	\end{frame}
	
	% ============================================================
	% SLIDE 52: REFERENCES - HARD REGULATORY SOURCES
	% ============================================================
	\begin{frame}{Regulatory References}
		\begin{thebibliography}{99}
			\bibitem{1} Financial Action Task Force. (2023). \textit{40 Recommendations of the FATF}.
			\bibitem{2} ACAMS. (2026). \textit{CAMS Study Guide}, 7th Edition.
			\bibitem{3} Digital Operational Resilience Act (DORA). (2022). Regulation (EU) 2022/2554.
			\bibitem{4} Wolfsberg Group. (2024). \textit{AML Principles for Correspondent Banking}.
			\bibitem{5} EU Directive 2018/843 (5th AML Directive). (2018).
			\bibitem{6} EU Directive 2019/1937 (6th AML Directive). (2019).
			\bibitem{7} FinCEN. (2023). \textit{BSA Regulations and Guidance}.
			\bibitem{8} Monetary Authority of Singapore. (2024). \textit{AML/CFT Notice 626}.
		\end{thebibliography}
	\end{frame}
	
	% ============================================================
	% SLIDE 53: FINAL QUIZ - TOUGH QUESTIONS
	% ============================================================
	\begin{frame}{Final Quiz (Hard Questions)}
		\begin{enumerate}
			\item What is the primary purpose of fuzzy logic in PEP screening?
			\item Which FATF Immediate Outcomes do supervisors prioritize for evaluating banks' preventive measures?
			\item What is the best action for a TCSP firm when a client requests confidentiality and is linked to high-risk industries?
			\item Under DORA, what is the key regulatory expectation for cross-border EU banks?
			\item What is the difference between FATF's "black list" and "grey list"?
			\item Why is "passive liveness detection" preferred over active methods in some contexts?
			\item What is the core vulnerability of private banking in the AML context?
			\item What action should a PSP take when entering a jurisdiction with an immature AML framework?
		\end{enumerate}
	\end{frame}
	
	% ============================================================
	% SLIDE 54: CONCLUSION
	% ============================================================
	\begin{frame}
		\begin{center}
			\Huge Thank You
			\vspace{0.5cm}
			\Large \textbf{Good Luck on Your CAMS Exam!}
			\vspace{1cm}
			\normalsize \textit{Remember: Risk-Based, Data-Driven, and Expert-Validated.}
		\end{center}
	\end{frame}
	
\end{document}